\documentclass[11pt, a4paper]{report}
\usepackage{amsmath}
\usepackage{amssymb}
\usepackage{graphicx}
\usepackage{epstopdf}
\usepackage{hyperref}
\usepackage[utf8]{inputenc}
\usepackage[T1]{fontenc}
\usepackage{geometry}
\geometry{margin=1in}
\usepackage{setspace}
\onehalfspacing
\usepackage{csquotes}
\usepackage{lipsum}

\title{The Adaptive Imperative}
\author{Albert Jan van Hoek}
\date{\today}

\begin{document}

\maketitle

\section*{Anti-capture, expendable parts, and the relocation of correction across scales}

We admire systems that can change themselves. We praise evolution because it adapts, science because it corrects itself, democracy because it can replace its leaders. Yet when a single person doubts themselves, we are inclined to call it weakness. The asymmetry is the first clue that we are looking at one phenomenon and miscounting it as several.

But the claim I most want to make is not that adaptation matters. That idea is old; it runs through Darwin, through cybernetics, through Ashby and Friston and complex-systems theory, and I will not pretend to have discovered it. The claim I want to make is narrower and, I think, new. It is that \textit{the force which keeps a system adaptable relocates as the system's agency relocates.} In a single mind it sits inside, as self-doubt. In a relationship it becomes reciprocal. In a society it becomes institutional. In an ecosystem it becomes ecological. The corrective moves outward as the number of agents grows --- and where it must sit turns out to be forced by the structure of who could possibly administer it, not chosen. A mind, moreover, is the rare kind of system that can install this force \textit{deliberately}, ahead of the collapse that would otherwise install it for free. That relocation, and that possibility of foresight, are the subject of this essay. The rest --- a substrate-independent principle of persistence, and the demotion of natural selection to one case of it --- is the scaffolding I need in order to state them precisely.

\section{A constraint above selection}

Natural selection is usually treated as the engine of biological adaptation: variation arises, the environment selects, the successful variants are retained, and the cycle turns again. There is a long tradition --- universal Darwinism, associated with Dawkins, Hull, and Dennett --- that tries to export this wholesale, treating ideas, technologies, and institutions as further populations of replicators undergoing selection. That export has drawn a sharp and, I think, partly correct objection, most forcefully from Lewontin: culture rarely contains anything like faithful replicators with measurable fitness differentials, and asserting that it does buys an appearance of unity at the price of content.

I want to make a different move, one that sidesteps that fight rather than rejoining it. Instead of claiming that minds and societies are \textit{also} Darwinian, I want to demote Darwinian selection from genus to species. Above it sits a more general constraint, which I will call \textbf{adaptive persistence}:

\begin{quote}
\textit{When the changes a system faces outrun what its fixed repertoire can handle, it can only persist if it can reconfigure itself.}
\end{quote}

Natural selection is then biology's particular way of meeting this constraint. Belief revision is cognition's way. Falsification, elections, and bankruptcy are some of the ways open to institutions. None of these needs to share a mechanism with the others. They need only share a function: the capacity to change the system's own structure in response to a world that has exceeded its current form. On this view you do not need replicators in culture, because culture is not required to \textit{be} selection. It is required only to instantiate the more general function that selection, in biology, happens to instantiate through replicators.

It will help to pin a few words before leaning on them, because they are kin but not synonyms. A system's \textit{adaptive space} is the set of configurations it can still reach from where it stands. \textit{Persistence} is its continued existence; \textit{viability} is persistence when that space keeps shifting underneath it, so that holding still is not on offer. \textit{Corrigibility} is the capacity to leave a configuration that has stopped working --- to move within the adaptive space rather than be stranded at a point in it. These form a chain rather than a list: viability requires corrigibility, and corrigibility requires that the adaptive space stay open. A fifth term, \textit{anti-capture} --- whatever keeps any single part from closing that space for the whole --- is the subject of the sections that follow, and the hinge the rest of the argument turns on.

This is not new ground so much as ground whose earlier surveyors deserve naming. Ashby's law of requisite variety holds that only variety can absorb variety --- that a regulator can control a system only if it has at least as much flexibility as the disturbances it faces. Friston's free energy principle offers a substrate-agnostic account of why anything that persists must behave as if it modelled, and corrected its model of, the world it lives in. And Hazen and colleagues have recently proposed a law of increasing functional information meant to sit above biological evolution, spanning the mineral, the biological, and the cultural. I am not claiming to displace any of these. I am borrowing their shared insight --- that persistence under change is substrate-independent --- and adding to it a structure they do not develop.

\section{Why this is not merely true by definition}

A reasonable reader will worry that ``systems that persist are the ones able to change themselves'' is true by definition and therefore says nothing. It is the same emptiness Lewontin warned of, and it has to be answered before anything is built on it.

The answer is that the constraint has content only under a specific condition, and is silent otherwise. A diamond persists for a very long time without any capacity to reconfigure itself, because its environment never asks it to: nothing in its surroundings outruns what its fixed structure can absorb. A species in a shifting climate has no such luxury. The principle therefore does not say ``everything that persists can adapt.'' It says something narrower and falsifiable: \textit{where no static configuration suffices, only reconfigurable systems persist.} The work is all in the antecedent. Specifying when the antecedent holds --- when change genuinely exceeds the fixed repertoire --- is what keeps the principle from collapsing into tautology, and it tells you where a counterexample would live: a system that faces such change, cannot reconfigure, and persists anyway would falsify the claim. I do not think you will find one. But you could look, which is what makes it a claim rather than a definition.

A fair worry is that the antecedent is only knowable after the fact --- that we read ``the environment outran the repertoire'' off the corpse and ``it did not'' off the survivor, smuggling the tautology back in. But the two quantities can be measured independently of the outcome, and ahead of it. Environmental volatility --- the rate and amplitude at which the relevant conditions change --- is estimable on its own; so is a system's rate of reconfiguration, how fast and how far it can move through its adaptive space. The antecedent holds when the first outpaces the second, and that is a prospective comparison, not a retrospective one. Population biology already runs it, under the name of evolutionary rescue: a population persists under deteriorating conditions only if its rate of adaptation outstrips the rate of deterioration, and the critical rate can be estimated before the population lives or dies. A system measured as flexible enough that nonetheless collapses, or as hopelessly outpaced that nonetheless persists unchanged, would falsify the claim. The condition is operational, not merely something we read off the ending.

\section{Anti-capture}

So far this is an account of how a system stays viable from the inside. The structure I am after appears once you ask a second question: not how a system changes itself, but what stops any single part of it from capturing the whole and shutting the changing down.

Consider the levels in turn. In a single mind, the danger is that one belief becomes unquestionable and the model stops updating. In a relationship of two, the danger is that one person's account of things dominates the shared account. In a society, the danger is that one firm, one fortune, one ideology grows large enough to foreclose the alternatives. In an ecosystem, the danger is that one species occupies the whole of the available niche space. At every level there is something that works against this: self-doubt, mutual revision, the institutional turnover of power and the openness of claims to refutation, predation and disturbance. Call the common job \textbf{anti-capture}. Its task is to prevent any single configuration from reaching the dominance that would destroy the variation the system feeds on.

It matters that this is not, at root, adversarial --- and this is why ``anti-capture'' is the right name and ``counter-force'' the wrong one. Self-doubt is not a mind at war with itself; science is not at war with itself. Predation and antitrust look oppositional, but the commonality across all the cases is not opposition. It is the prevention of capture. And anti-capture is not a separate phenomenon from corrigibility; it is corrigibility distributed --- the same holding-open that a mind performs on itself, now performed by the parts on one another. Self-doubt is a mind holding its own dominant belief open to revision. Mutual revision is two minds holding each other's open. Redistribution is a society holding its own winners open to replacement. The same act recurs; only the number of agents changes.

This is not a bare analogy. Its sharpest empirical form is old and well established. In 1966 Robert Paine removed a single predatory starfish from a stretch of rocky shore and watched the consequence: freed from predation, one mussel species expanded and crowded the others out, and the diversity of the whole community roughly halved. The lesson generalises with unusual cleanness --- removing the thing that holds the dominant competitor in check does not free the system, it impoverishes it. Around this sit findings that fill in the picture: niche construction, in which organisms build the very niches they then occupy (the advantage is one you make, not only one you find); Red Queen dynamics, in which each advance provokes a counter-advance; and the coexistence theory that explains why the strongest competitor so often does \textit{not}, in fact, take all. The idea of intermediate disturbance --- that moderate, recurring disruption maximises diversity --- is the closest ecological echo of the argument here, though it is genuinely contested in its strong form, so I lean on Paine, who is not.

\section{What survives is the process, not the parts}

Push on anti-capture and it turns into something stronger than the prevention of dominance. To prevent any part from capturing the whole is to prevent any part from becoming permanent --- and permanence is the real threat that dominance is only the visible form of. A dominant belief is one that has escaped revision. A monopolist is a firm that has escaped replacement. An entrenched leader is power that has escaped turnover. Anti-capture, stated at full strength, is the keeping-mortal of the parts.

Which means every adaptive system --- under the antecedent condition, where no static configuration suffices --- requires not only retention but turnover. Evolution preserves genes and destroys them. Brains preserve beliefs and destroy them. Science preserves theories and destroys them. Relationships preserve a shared narrative and revise it. Civilizations preserve institutions and dismantle them. What persists across all of this is not the components. It is the process by which components are allowed to die. Adaptive systems do not survive because their parts are immortal; they survive because their parts are expendable, while the process that replaces them is kept open.

This is almost paradoxical, and it is the deepest connection to the intuition the essay began from. The sacred thing is not the gene, the belief, the leader, or the institution. It is the openness of the process by which each of these is permitted to go. The objection to calling self-doubt weakness lands here exactly: self-doubt is not the doubting of a self. It is the willingness to let a belief die in order to keep the capacity to learn. A mind that can do this performs, within a single lifetime, what evolution performs across generations --- it lets the model die instead of dying with the model.

\section{Layered persistence}

Said flatly, ``the parts are expendable'' is false, and a reader is right to balk. Make the genetic machinery as disposable as a phenotype, or a constitution as revisable as a parking regulation, and you do not get adaptation; you get noise, and then collapse. The resolution is that expendability is not uniform across a system --- it is stratified. Every adaptive system that lasts is layered, with a slow, heavily shielded core and a fast, freely revised edge, and the shielding of the core is precisely what \textit{funds} the fluidity of the edge. You can let phenotypes die in every generation because the genetic code that builds them is conserved across hundreds of millions of years. You can let policies turn over because the constitutional frame that lets policy be made and unmade is held nearly still. The expendability of the surface is bought by the rigidity of the base. They are not in tension; they are complements. So the earlier claim has to be read with its layer attached: parts are expendable \textit{where adaptation happens}, and that layer stays open only because a deeper one does not.

This has good names. Simon argued that the complex systems which survive are nearly decomposable --- built from stable sub-assemblies whose internal couplings are strong and fast and whose external couplings are weak and slow --- and his parable of the two watchmakers is the point in miniature: the one who builds in stable sub-units can be interrupted and resume; the one who must hold every part in place at once is never finished. Developmental biology calls the shielding of core outcomes \textit{canalisation}, after Waddington --- the buffering that keeps the same body plan emerging across a wide range of genetic and environmental perturbation. Engineers rediscovered the same shape as the hourglass: the internet's narrow, frozen waist at the protocol layer, which by \textit{ending} diversity at that one level let it explode at every level above and below. In each case a layer is deliberately taken out of the adaptive game so that the game can be played faster everywhere else.

This also dissolves what at first looks like a counterexample to the whole argument: the cases where capture does not exhaust a substrate but founds a new one. A cell engulfs another and does not digest it; the captive becomes the mitochondrion, a permanent internal monopoly --- and that very entrenchment is what powered all complex life. Free-living cells surrender their autonomy to a multicellular body, and the body becomes the unit that now adapts. These are the major transitions in evolution, and they share a shape: independent parts are captured into a higher unit, the capture is frozen into that unit's core, and a fresh adaptive space opens one level up. This is not a breach of anti-capture; it is the same thing read along the vertical. Terminal capture closes the adaptive space and opens nothing above it --- the monopoly that ends up owning a dead market. Scaffolding capture freezes a layer and opens a new space above it. The test of any capture is therefore not ``did a part come to dominate?'' but ``did the adaptive space close, or relocate upward?''

Two cautions keep this from becoming a licence for permanence. First, no core is immortal; it is slow, not exempt. The genetic machinery does change, across deep time, and constitutions are amended and occasionally replaced. A core is a layer whose correction runs on a long clock --- not a layer outside the process. Second, a frozen core that scaffolded one era can become the constraint of the next: the lock-in, the standard that no longer fits but that everything is now built upon, the foundational law that has outlived the world it was written for. When the layers above can no longer route around the core, the scaffolding has turned terminal, and the slow clock has to move after all. Which is why the magnitude question of the next sections is really one question asked once per layer --- how fast should \textit{this} layer correct --- and why the architecture of a durable system is the whole schedule of those rates: fast at the surface, slow at the base, and none of them stopped.

\section{What anti-capture is doing}

Why does the danger of capture arise at all? Because the unit that gets optimised is rarely the unit whose viability is at stake. Selection acts locally; viability is global. A cancer cell is selected, locally, to divide, and the organism dies. A firm is rewarded, locally, for taking the whole market, and the market dies. A nation gains, locally, by defecting, and the alliance dies. A belief is reinforced, locally, by being comforting and identity-confirming, and the mind that holds it stops learning. None of these is best described as free-riding; each is local optimisation closing global adaptability. That is the problem anti-capture exists to meet, in its most general statement: it is whatever keeps local optimisation from destroying global adaptability. And the cheapest instance is the last one --- a single mind --- which is why self-doubt belongs to the same family as predation and antitrust, and why this essay could begin where it did.

This is not the environment ``stepping in to preserve diversity,'' as though preservation were a purpose; ecosystems have none. The mechanism is blind and causal: the locally optimising part consumes the very substrate that produced its advantage, and the larger system loses the variety it was running on. What this adds to the picture is timing, and an order of harm that is easy to misstate. When one part takes everything, it is everyone else's loss first; the part's own loss comes later, on a delay, when the substrate collapses beneath the position it worked so hard to secure --- the predator that has eaten all its prey, the monopoly that owns a dead market. Dominance is not immediately self-defeating. It is \textit{eventually} self-defeating, through a mechanism --- substrate exhaustion --- that can be named and, in principle, measured.

There is a human version with a precise name. People often resist a ceiling on great wealth not because they hold such wealth, but because they expect, or hope, to. Economists call this the prospect of upward mobility --- Bénabou and Ok's account of why some who would benefit from redistribution nonetheless oppose it. The structure is a trap. By protecting the \textit{absence} of a ceiling, in order to keep their own imagined route to the top open, people make it possible for someone to reach a height that closes the route for everyone, themselves included. The very openness they are trying to preserve is destroyed by the dominance that the lack of a ceiling permits.

\section{The upgrade: substrates that can install anti-capture themselves}

Here is the difference that matters most, and it is a difference between substrates.

A species cannot read the theory of evolution and mutate more wisely. DNA carries no representation of the population-level process it is subject to; anti-capture, for a species, is always imposed from outside, after the fact, as predation or famine or extinction. A mind is different. A mind can hold a model that includes its own dynamics and act on that model. It can install the corrective \textit{before} the catastrophe that would otherwise install it --- it can doubt the belief before reality refutes it, choose the restraint before the collapse compels it.

I want to be careful about what property is doing the work, because it is easy to reach for words that claim too much. It is not consciousness, and not ``agency'' in any deep metaphysical sense, and --- importantly --- not full comprehension of one's own substrate. I do not fully understand how DNA computes, or how a neural network does, and I do not need to in order to draw the relevant line. The property that matters is narrower: \textbf{representational closure under self-reference} --- whether a system can build a model that includes its own behaviour and let that model feed back into what it does. DNA cannot. A brain can. An artificial network plausibly can, which is exactly why stating the principle in substrate-neutral terms earns its keep rather than being a flourish. The property comes in degrees: an institution holds a partial model of itself --- a constitution is a society's imperfect self-model --- and so sits somewhere in between. The line I need is only the coarse one: some substrates can host a self-model that changes their behaviour, and some cannot.

Within a single system, the imperative that follows is not yet an ethical claim but a conditional one. If survival and the capacity to learn are valued, corrigibility must be preserved --- as breathing requires oxygen. Nothing more normative than that conditional is needed here, and a system that goes on existing through change already behaves as though it holds the value. The complication arrives only when there is more than one agent, which is where the relocation finally bites.

\section{Where correction has to sit}

Inside a single mind, anti-capture is reflexive: the mind is its own corrective, because self-doubt is a mind holding itself open, and a mind has the self-reference to do this. Across competing parties it is tempting to say the corrective simply \textit{cannot} be self-administered by the dominant one --- that asking the position whose excess the corrective exists to check to administer that check against itself is incoherent. But that is too strong, and history refuses it. Dominant parties do sometimes limit themselves, and when they do it is the upgrade at work, not a breach of it: Bismarck built the first welfare state to forestall a revolution he could see coming, and firms standardise and disclose to pre-empt the antitrust or the backlash that would cost them more. A dominant party that genuinely foresees the substrate collapsing under it, and that will still be standing when it does, has every reason to install the corrective on itself.

The real asymmetry is not that elites cannot self-model; it is that dominance usually lets them not have to. The thing the local-versus-global section already named --- that capture lands as everyone else's loss first and the winner's loss only on a delay --- is exactly what blunts elite foresight. A dominant party can often externalise the collapse it foresees, or defer it past its own tenure, so that its self-model correctly reports that \textit{it} will not be the one to pay; and then self-limitation is not rational. It becomes rational only when that escape is closed --- when the party can neither offload the cost nor outrun it, as Bismarck could not outrun a domestic revolution. So the corrective is reliably installed by whoever bears the collapse, and dominance is, among other things, the power to redistribute who bears it. That is why, as a standing matter, the reliable site of the corrective is the rest --- the parties who cannot externalise the loss, because the substrate is theirs. Not because elite restraint is impossible, but because it is the exception that needs the escape route closed first, while for the rest the incentive is already built in.

This is also where the conditional logic of the previous section stops transferring cleanly, and where a genuine value has to be named rather than assumed. It is one thing to say a system's persistence requires its corrigibility; it is another to say \textit{we should build societies that protect adaptive openness}. The second selects the persistence of the shared space as a goal --- and at the collective scale that goal is contested, because the dominant party's revealed preference is for its own permanence, not for the openness of the commons. The gap between ``each agent acts to persist'' and ``the shared space stays open'' is the same local-versus-global problem from earlier, now in its multi-agent form --- the form economists call free-riding --- and it does not close itself. This is the one place an honest argument must mark a commitment: that we value keeping the adaptive space open. Given that value the rest follows; without it the descriptions remain descriptions. The earlier conditional was logic; this is a choice, and pretending otherwise would be the sleight of hand.

With the value named, the political structure becomes clear, and it reframes what first looks like a demand for individual sacrifice. The rule is not best understood as ``the dominant should restrain themselves.'' It is a rule, endorsed by the many, that binds whatever party occupies the dominant position --- including the future, hypothetical version of oneself who might one day arrive there. The cost falls on a position, not a person; on a self one might become, while the benefit of an open system accrues to the self one already is. That asymmetry is what makes such a rule endorsable where self-restraint is not. You are not asked to surrender what you hold. You are asked to agree that the height you aspire to should carry a ceiling --- and since you would want that ceiling honoured if you arrived, agreeing to it is not generosity but consistency. From the inside it can feel like agreeing to be taxed if you grow rich; in structure it is the field imposing a limit on whoever stands highest. Both are true, and the distance between them is the seam being honest about itself.

This is the shape of contractualist arguments made behind a veil of ignorance --- Rawls's device of choosing the rules without knowing which position you will occupy --- but the version that arises here is sturdier than the thought experiment, because it does not require anyone to pretend. The person who wants to become rich genuinely does not know whether he will. The uncertainty that makes the rule fair is a real feature of his situation, not a hypothetical he has to talk himself into. The same shape is visible in how human groups actually sustain cooperation: enforcement is reliably placed outside the contestants --- in third parties, in institutions, in the commons-governance arrangements Ostrom documented, where communities install on themselves rules no single member could be trusted to apply to their own advantage. The contribution here is the claim that this placement is not a contingent fact about human politics but follows from where the cost of collapse falls: the corrective settles, reliably, on whoever cannot escape that cost --- which at the collective scale is the rest, and almost never the party the corrective limits.

A political objection answers itself inside this structure. If the rule binds the top, what stops the top from blocking it? The answer is in the numbers: the rule is wanted by everyone not currently at the top, which under any real distribution is almost everyone. What erodes that majority is precisely the prospect-of-upward-mobility belief, which persuades the aspirant to identify with the position he hopes to reach rather than the one he occupies.

What I have just described in the language of wealth and ceilings is only one instance, and it is worth keeping the theory above the example. Abstractly, societies persist by developing mechanisms that prevent durable capture of their adaptive space. A ceiling on accumulation is one such mechanism. Competitive elections are another, and term limits another, and antitrust another, and decentralization another, and the openness of scientific claims to peer refutation another. The theory does not depend on any of them in particular; it survives the replacement of any specific arrangement by a better one. I develop the wealth case because it makes the relocation of the corrective concrete --- it shows clearly why the limit reliably falls to those who would bear the loss rather than to the party it limits --- not because the argument stands or falls with that policy. And none of these mechanisms is a destination the system is climbing toward. Evolution has no arrow; Gould spent a career insisting, correctly, that it trends toward no goal, and the commons in the sense meant here is not a goal either. It is a structural condition: a domain in which no single party is permitted to capture the substrate that all draw on.

\section{How much}

Agreeing that a corrective is necessary, and that it reliably falls to the rest, leaves the largest practical question open: how strong should it be? A ceiling on accumulation could be set high or low; predation can be light or heavy; doubt can be a whisper or a paralysis. The mechanism does not fix the magnitude, and the magnitude is not a detail to be settled afterward but part of the theory, because the dynamics of the whole system depend on it.

The reason is that anti-capture fails in both directions. Too weak, and it does not prevent capture: the dominant party monopolises, the substrate is exhausted, the diversity collapses --- the starfish removed, the mussels taking the shore. Too strong, and it removes the very gradient the system runs on, or destabilises it into noise. Biology supplies the cleanest illustration of the second failure. A population needs a mutation rate high enough to generate the variation selection acts on; but past a threshold --- Eigen's error catastrophe --- mutation becomes so frequent that no adaptation can be stably inherited, and the population's information dissolves faster than selection can preserve it. There is too little variation, and there is too much, and viability lives in the band between. The same shape recurs at every scale. Too little self-doubt is dogmatism; too much is an inability to hold any belief long enough to act on it. Too little mutual revision lets one partner dominate; too much means no shared understanding is ever stable enough to build on. Too little redistribution permits monopoly; too much can flatten the differences a system also draws on.

I want to be precise about what this gives us, because it is easy to over-read. It does \textit{not} give an optimum in the sense of a single correct number waiting to be discovered. ``Optimal for what?'' is a fair and unavoidable question with no purely factual answer, because the things a corrective governs --- the rate at which a system changes, the complexity it can sustain, the diversity it harbours, its stability over time --- trade off against one another. A setting that maximises diversity may cost stability; one that maximises the speed of adaptation may cost the complexity that only accumulates when structures persist. There is, at best, a frontier of viable settings, and the choice of where on that frontier to sit is itself a value decision --- exactly the kind that should be left to those who will live with it, and revisited as conditions change, rather than fixed once. The science can map the frontier; it cannot choose the point.

What the theory contributes is the recognition: that the strength of the corrective is a parameter --- and, as the layered account showed, one parameter per layer rather than one for the whole system --- that it governs the trade-off among speed, complexity, diversity, and stability, and that there is a band outside which the system fails in one direction or the other. Where within that band a generation chooses to live is properly open, and characterising --- for systems of different kinds --- the shape of that frontier and the location of its edges is among the more interesting questions left for empirical work. The claim here is only that the question belongs inside the framework rather than after it.

\section{A ratchet of emergence}

Four pressures on the argument turn out to be one, and naming it is the last structural piece. Each asks, in a different register, the same question: where does the corrective actually come from, and what forces it? The answer is a single loop, and it is the engine under everything above. A false or dominant configuration has slack --- room to persist out of true --- but the slack is finite, and as it runs out reality begins to press. That pressure is what I have been calling tension, and it does two things at once: it is the cost arriving, and it is a gradient. Along the gradient variation explores, blindly and without foresight, and the configurations that relieve the pressure are selected and retained, ratcheted in as the new ground the next round runs from. Reality biting and a solution emerging are not two events; they are one event seen from its two ends --- depletion of slack on one side, generation of a successor on the other.

How long the slack lasts is not fixed; it is set by how tightly a configuration bears on a substrate that cannot be renegotiated. A false aerodynamic model has almost none --- the wing cannot be argued out of falling, and the bite is immediate. A false economic model has a great deal, because the social surface can be reshaped to fit it for a long time. A pure convention, answering to no physical floor at all, has more still. But none of the slack is infinite, because the negotiable layers rest in the end on floors --- energy, food, bodies, a climate --- that no configuration can coerce, and against those a false model quietly accrues fragility until they give. The clock differs by substrate; that there is a clock does not.

This is also the answer to the hardest of the four pressures: if capture mostly eats its substrate, how does it ever become the rigid core that scaffolds a level above? The monopoly does not convert --- it is conquered from above. Most captures are terminal; the system dies, or falls back a level. Scaffolding is the selected exception, and its mechanism is exactly anti-capture installed one rung up. Under the tension of a capture, the configurations that get retained are the rare ones in which a higher-level unit comes together that can control the captor --- disarm its extraction and enrol its dominance as a function. The mitochondrion did not volunteer to stop eating; the eukaryotic cell took its autonomy, most of its genome migrating to the nucleus so that it could no longer act as an independent replicator. A multicellular body is viable only once it installs apoptosis and immune surveillance against the cells that defect --- which is to say, against cancer. So terminal capture becomes scaffolding precisely when a higher level can install anti-capture on the captor, and stays terminal when no such level forms: the cancer that outruns surveillance kills the host, the firm that captures its regulator kills the market. That is a forward condition, not a story told backward; we are not saying captures scaffold because they happened to survive, but naming what has to be true for one to. And the coming-together is the constructive face of the same event whose regulatory face is the control. Emergence here is a network closing into a new unit --- which is why, for the systems that manage it, the word carries something cooperative and not only corrective.

Read through the same loop, the collective-action objection stops being a refutation and becomes a fork. Bearing the cost of collapse is an incentive, not a capacity, and a million people each losing a little will not reliably organise against the one who gains a lot --- diffuse costs and concentrated benefits paralyse, as Olson showed. But finitude still guarantees that the corrective arrives; the runway ends. The only question is by which route. If the rest can pay the cost of coordinating, they install the corrective deliberately, ahead of collapse --- the upgrade. If they cannot, they drop to the branch the essay has already named, where collapse installs it for them, at full price. Paralysis does not show that no correction comes; it shows which one does, the expensive instead of the cheap. And the thing the objection calls missing --- a mechanism to lower the cost of organising the rest --- is not missing; it is anti-capture itself. Institutions, shared models, common knowledge are precisely the machinery that makes the deliberate branch reachable. An argument like this one, if it works, is a small instrument of that kind.

The last pressure marks the honest limit of the whole apparatus. The rate-comparison that operationalised the antecedent --- volatility against reconfiguration speed --- works only where the change is continuous enough to be measured in advance. Genuine novelty, the shock undefined until it lands, is by definition the case where no such measurement exists, and the framework should promise no foresight there, because there is none. What it promises instead is the other half of the ratchet: where the specific shock cannot be prepared for, what survives is retained generativity --- diversity, redundancy, reserve slack --- not because it is aimed at the shock but because it keeps the capacity to generate a response to whatever comes and lets the shock select. This is the deep reason the frontier of the previous section should err, under real uncertainty, toward more slack than point-optimisation would advise: the slack is not waste, it is the material emergence draws on when the future arrives unannounced. It is also the structural form of the ending this essay is moving toward --- you cannot foresee what will refute you, so you keep the capacity to be refuted.

One discipline keeps all of this from collapsing into the emptiness it was built to avoid. \textit{Emergence} explains nothing if it means only ``and then a solution appeared''; ``the tension produced whatever resolved the tension'' is true by definition, and Lewontin would be owed his due. The word earns its place only in the weak sense --- a resolution derivable in principle but only by running the process, never readable in advance from the parts --- and only when it is cashed out as the loop itself: variation generated under the tension, which shapes the search by fixing what would count as relief; feedback selecting what actually relieves it; retention ratcheting it in. That is not an alternative to selection. It is selection's generative front end, working where the answer cannot be foreseen --- evolution running in the dark. The framework around evolution is its mechanism, and emergence is the name for what that mechanism is doing at the moment the next configuration is found.

\section{What we do not get, even at the end}

It is tempting, having built all this, to reach for a final claim about truth --- to say that truth simply \textit{is} the long-run survivor of correction, the model left standing after enough refutation. The temptation is strong because it would unify the epistemology with everything above it: reality exists, models compete, the less corrigible die, and what remains is true. It is also the one claim in the vicinity I think is wrong, and seeing exactly why is worth more than the unification it would buy.

Selection tracks viability, not truth --- and the two are not the same thing, even when they travel together. A model can survive a great deal of correction without being the last word, and Newtonian mechanics is the instructive case, because what happened to it was not simple falsification but something this paper's own vocabulary anticipates. Its picture of the world --- absolute time, force acting at a distance --- was superseded; its predictive structure was kept as an approximation whose domain we learned to bound. The claims were expendable; the working part survived. So two centuries of success certified its \textit{viability} --- its fitness across a range of circumstances --- and not its final accuracy about what the world is. A scientific realist will object, rightly, that sustained success on novel predictions is \textit{some} evidence of approximate truth; and so it is, defeasibly. But evidence is not a certificate, because what the process optimises is viability, and a less accurate model that is cheaper or better-fitted can outlast a truer one. Survival is the test the process actually runs. Truth is not guaranteed to be the thing that passes it.

The pragmatist temptation, Peirce's, is to close this gap by definition --- to say that truth simply \textit{is} the end of inquiry, the convergence of correction. It inherits a familiar difficulty: it cannot say what inquiry converges \textit{toward} without either pointing past the process to a reality the process only approximates, or sliding into the relativism it meant to avoid. Saying ``not in a relativistic sense'' does not earn the exemption. Something outside the correction has to anchor it --- and that something is exactly what the process cannot supply from inside.

There is a real pull toward truth inside this, though, and it is worth being exact about its shape, because a critic is right that viability and truth are not equally far apart at every scale. The gap between them is a function of how deeply a system's viability depends on getting reality right, and that depth grows with the system's reach. A tribe can persist for centuries on a false cosmology --- lightning as an angry god --- because nothing it does puts that model under load; its tolerances are wide. A civilisation that etches circuits a few atoms across cannot, because the intervention simply fails unless the quantum mechanics behind it is very nearly true. As manipulative reach deepens, the tolerances tighten, and viability does begin to force a convergence toward truth. That is the grain of sense in the pragmatist's hope, and it should be granted.

But it is local, not global, and granting it whole would be the mistake. Viability constrains a model only along the dimensions the system actually probes. The same civilisation that prints a flawless transistor can carry a wildly false account of its own social dynamics, or of its long-run dependence on a climate, and be killed by precisely the dimension its capability never put under load --- exact where it intervenes, blind where it does not. And there is a darker turn, which is anti-capture reappearing at the level of truth itself: capability is also the power to keep a false model viable \textit{by reshaping the world to fit it} rather than the model to fit the world --- to suppress the feedback that would correct it, externalise its costs, manufacture the consent that holds it in place. A system powerful enough can keep a false model not because it has converged on truth but because it has captured its own epistemic environment, made itself too strong to be corrected. But the evasion is bounded, for the reason the engine gave: capture buys wiggle on the layers that can be renegotiated and none on the floor beneath them. You can manufacture consent about the economy for a generation; you cannot manufacture the harvest that feeds it. The capture is real, and it is long, and it is never permanent, because some substrate underneath was never up for negotiation. So capability cuts both ways. It tightens the tolerances along the axes it probes and loosens them, by evasion, along the axes it would rather not --- which is only to say that the wish to be finished correcting grows more dangerous, not less, as the power to enforce it grows.

So here is the more exact landing, and I think the stronger one. The model that survives correction is the best we have, and its survival is evidence against our \textit{certainty} rather than against its truth --- but it is no warrant for finality either. And the urge to treat it as final --- to declare the inquiry closed, the survivor certified, the doubt no longer required --- is \textit{itself} a bid for capture. It is a claim reaching for the one thing anti-capture denies every part: permanence, immunity from replacement, escape from the process that produced it. The principle of this paper therefore turns on the paper. The wish to stop, to possess the answer, to be finished correcting, is the cognitive capture the corrective must keep dissolving --- in a single mind first, where it is the oldest form of the temptation, and at every scale above.

This is the deepest connection to where the essay began. Self-doubt was never really about doubting yourself. It is the willingness to let a belief die in order to keep the capacity to learn --- the smallest instance of the largest pattern: that adaptive systems persist not because their parts are immortal but because their parts are expendable, while the process that replaces them is kept open. Genes are expendable. Beliefs are expendable. Leaders, theories, institutions, even the convictions in this paragraph, are expendable. The only thing it is rational to treat as something close to sacred is the openness of the process by which the expendable is allowed to go. That word has to be qualified, though, because the regress is real and worth naming: openness is not sacred in itself. It is what becomes worth protecting once persistence and the capacity to learn are valued, and that valuing is the one premise the whole structure rests on and cannot itself supply. Press far enough --- why protect openness? because it preserves adaptation; why adaptation? because we wish to persist and to learn --- and the chain ends not in a fact but in a commitment, which the argument can build on but cannot reach beneath.

This framework is no exception to itself, and it would be incoherent if it tried to be. It is a configuration in the same adaptive space as everything it describes. If it is wrong, it should be replaced; if it is partial, it should be expanded; if it is useful, it has earned no more than provisional survival, for as long as it keeps passing the test. The one move it is not entitled to is the move it warns against: to exempt itself from correction, to become a part that has escaped revision. A theory immune to its own principle would be the exact form of capture the principle names. So this one holds itself to its own terms.

A mind is the kind of system that can see this and act on it: that can keep the process open deliberately, install the corrective before collapse imposes it, and resist --- in itself, and in the arrangements it builds with others --- the standing temptation to declare the correcting finished. Survival has never belonged to the strongest, nor to the most certain. It belongs to whatever keeps itself able to change, and, where there are many, to whatever can agree to hold open the space in which the changing happens, at a strength it leaves its successors free to revise. Even this conclusion asks to be doubted. That is not a weakness of the argument. It is the argument.

\section*{On precursors}

This essay stands on work it does not try to summarise. The idea that persistence under change is substrate-independent is anticipated by Ashby's law of requisite variety, by Friston's free energy principle, and by the recently proposed law of increasing functional information of Hazen and colleagues; the attempt to generalise selection across domains, and the strongest objection to it, belong to the universal-Darwinism tradition (Dawkins, Hull, Dennett) and to Lewontin respectively. The ecological backbone is Paine's keystone-predation work, with niche construction, Red Queen dynamics, and coexistence theory alongside it, and the operational reading of the antecedent borrows the rate-comparison of evolutionary-rescue theory. The layered account --- a shielded slow core funding a fluid fast edge --- draws on Simon's near-decomposability and Waddington's canalisation, and the cases where capture founds a new level rather than exhausting a substrate are the major transitions in evolution of Maynard Smith and Szathmáry. The political structure draws on Bénabou and Ok's prospect-of-upward-mobility hypothesis, on the contractualist tradition associated with Rawls, on Ostrom's account of how communities govern a commons, and on Olson for why bearing a cost is not yet the capacity to act on it. The threshold argument borrows its sharpest case from Eigen's error catastrophe; the closing caution about truth is set against the pragmatist account associated with Peirce; the limit set by genuine novelty is the black-swan problem named by Taleb. The recurring tension between local and global optimisation has relatives in multilevel-selection theory and in the evolutionary account of cancer as a breakdown of cellular cooperation, and \textit{emergence} is used throughout in the weak sense associated with Bedau --- a resolution derivable only by running the process, not foreseeable from the parts. The contribution I have tried to add is not the general level itself but four things: anti-capture as corrigibility operating one scale up and, at full strength, the keeping-expendable of a system's parts; the stratification of that expendability into layers, so that a rigid core funds a fluid edge and a capture that freezes one layer can open the next; the relocation of the corrective from the reflexive to the collective, sited with whoever bears the collapse; and the strength of the corrective as a per-layer parameter the framework must own rather than defer. Beneath these is one mechanism, offered not as new but as the engine they are facets of: variation generated under tension, selected and ratcheted, with emergence its name for the generative step where the resolution cannot be foreseen.

\end{document}